\documentclass[10pt,a4paper]{amsart}
\usepackage[margin=19mm]{geometry}
\usepackage{amsmath,amssymb,mathtools}
\usepackage[scr=rsfs,cal=euler]{mathalfa}
\usepackage{xfrac}
\usepackage[unicode,hidelinks,bookmarks=false]{hyperref}
\newcommand{\ie}{i.e.\ }
\newcommand{\cf}{cf.\ }
\newcommand{\resp}{respectively\ }
\newcommand{\Subsection}{\subsection}
\providecommand{\nobreakdash}{\nobreak\hbox{-}}
\setlength{\emergencystretch}{2em}
\multlinegap0pt
\usepackage{amsopn}
\usepackage{amsaddr}
\usepackage{xy}
\usepackage{amscd}
\usepackage{color}
\usepackage{verbatim}
\usepackage{latexsym}
\setlength{\topmargin}{0mm} \setlength{\headheight}{0mm}
\setlength{\headsep}{10mm} \setlength{\textheight}{220mm}
\setlength{\footskip}{15mm} \setlength{\marginparpush}{20pt}
\setlength{\oddsidemargin}{5mm} \setlength{\evensidemargin}{0mm}
\setlength{\textwidth}{150mm} \setlength{\marginparsep}{0mm}
\setlength{\marginparwidth}{20mm} \setlength{\parindent}{0em}
\setlength{\parsep}{20em}
 \newtheorem{theorem}{Theorem}[section]
\newtheorem{lemma}[theorem]{Lemma}
\newtheorem{proposition}[theorem]{Proposition}
\newtheorem{definition}[theorem]{Definition}
\newtheorem{corollary}[theorem]{Corollary}
\newtheorem{notation}[theorem]{Notation}
\newtheorem{example}[theorem]{Example}
\newtheorem{question}[theorem]{Question}
\newtheorem{remark}[theorem]{Remark}
\newenvironment{prf} {{\bf Proof.}}{\hfill $\Box$}
\newtheorem{thm}{Theorem}[section]
 \newtheorem{cor}[thm]{Corollary}
 \newtheorem{lem}[thm]{Lemma}
 \newtheorem{prop}[thm]{Proposition}
 \newtheorem{rem}[thm]{Remark}
 \newtheorem{defn}[thm]{Definition}
 \newtheorem{exam}[thm]{Example}
 \newtheorem{dis}[thm]{Discussion}
\newenvironment{emphit}{\begin{itemize}\em}{\end{itemize}}
\newcommand\RR{{\mathbb{R}}}
\newcommand\CC{{\mathbb{C}}}
\newcommand\QQ{\mathbb{ Q}}
\newcommand\NN{\mathbb{ N}}
\newcommand\ZZ{\mathbb{ Z}}
\newcommand\TT{\mathbb{ T}}
\newcommand\HH{\mathbb{ H}}
\newcommand\BB{\mathcal{B}}
\begin{document}
\title[Adaptive resolution frames: A multilevel framework in Hilber]{Adaptive resolution frames: A multilevel framework in Hilbert spaces}
\author{2020 Mathematics Subject Classification. 42C15. ęywords adaptive resolution frame, diagonal operator, level-wise reconstruction, stability, weighted reconstruction; Jahangir Cheshmavar}
\date{}
\maketitle
\noindent\emph{Source and licence.} Adaptive resolution frames: A multilevel framework in Hilbert spaces. Version arXiv:2608.18879v1 (CC BY 4.0). This material is distributed under the Creative Commons Attribution 4.0 International licence, http://creativecommons.org/licenses/by/4.0/, as recorded in the arXiv OAI licence metadata for this exact version and shown on the abstract page https://arxiv.org/abs/2608.18879v1. Full rights notice: licenses/arxiv-2608.18879-CC-BY-4.0.txt.\par\smallskip
\noindent\textit{Editorial scope.} Counted unit: the original \texttt{prop} environment \texttt{Prop.22} at source lines 500--539 together with its complete proof at lines 540--658; the delivered document keeps source lines 376--659 so that every referenced label and every citation resolves inside it. Native editable source is the paper's own concatenated TeX; the AMS wrapper is editorial formatting only. No publisher class, upstream script or unrelated artwork is redistributed.
\section{\textbf{The results}}
Throughout the section, let $\{f_i\}_{i\in I}\subset\mathcal H$ be
an $AR$-frame with frame bounds $A, B$, positive weights
$\{\lambda_i\}_{i\in I}$, partition $I=\bigcup_{k=1}^\infty I_k$,
and with orthogonal decomposition
\begin{eqnarray*}
\mathcal H=\bigoplus_{k=1}^{\infty}\mathcal H_k, \qquad \mathcal
H_k=\overline{\operatorname{span}}\{f_i:i\in I_k\},
\end{eqnarray*}
where each $\{f_i\}_{i\in I_k}$ is a frame for $\mathcal H_k$.
Also, let
\begin{eqnarray*}
S_k:\mathcal H_k\rightarrow\mathcal H_k, \,\ S_kx:=\sum_{i\in
I_k}\lambda_i\langle x,f_i\rangle f_i,
\end{eqnarray*}
for any $x\in\mathcal H_k$, be the level-wise frame operators, and
$S=\sum_{k=1}^\infty S_k$ be the full frame operator.\\

Our first result establishes a block operator representation for
adaptive resolution frames and provides a convergent iterative
reconstruction method with an explicit error estimate.

\begin{prop} \label{Prop.1}
Suppose that $T\in\mathcal B(\mathcal H)$ satisfies $T(\mathcal
H_k)\subseteq\mathcal H_{k+1}, \,\ k\ge1$, and
$\rho:=\|S^{-1}T\|<1$. Then the following statements hold:
\begin{itemize}
\item[(i)] There exists an orthonormal basis of $\mathcal H$ in
which every $S_k$ is diagonal, the frame operator $S$ is block
diagonal, and $T$ is strictly lower triangular with respect to the
decomposition $\mathcal H=\bigoplus_{k=1}^{\infty}\mathcal H_k$.

\item[(ii)] For every $x\in\mathcal H$, define the sequence
$\{u_n\}_{n\ge0}$ by $u_0=0$, and
\begin{eqnarray} \label{eqn-90}
u_{n+1} = u_n + S^{-1}\bigl(x-(S-T)u_n\bigr), \,\ n\ge0.
\end{eqnarray}
Then $\{u_n\}_{n\geq 0}$ converges in norm to the unique solution
$u=(S-T)^{-1}x$. Moreover, $\|u-u_n\| \le \rho^{\,n}\|u\|, \,\
n\ge0$.
\end{itemize}
\end{prop}

\begin{proof}
Since $\mathcal H=\bigoplus_{k=1}^{\infty}\mathcal H_k$ is an
orthogonal decomposition, every $\mathcal H_k$ is a closed
subspace of $\mathcal H$. For each $k$, the operator
\begin{eqnarray*}
S_kx=\sum_{i\in I_k}\lambda_i\langle x,f_i\rangle f_i
\end{eqnarray*}
maps $\mathcal H_k$ into itself. Moreover, if $x\in\mathcal H_j$
with $j\neq k$, then $\langle x,f_i\rangle=0, \,\ i\in I_k$,
because $\mathcal H_j\perp\mathcal H_k$. Hence $S_kx=0$.
Therefore, $S=\sum_{k=1}^{\infty}S_k=\bigoplus_{k=1}^{\infty}S_k$,
so the frame operator is block diagonal with respect to the
decomposition $\mathcal H=\bigoplus_{k=1}^{\infty}\mathcal H_k$.
Since every $S_k$ is bounded, positive, self-adjoint and
invertible on $\mathcal H_k$, the spectral theorem provides an
orthonormal basis $\{e_i^{(k)}\}_{i\geq 1}$ of $\mathcal H_k$
consisting of eigenvectors of $S_k$,
\begin{eqnarray*}
S_ke_i^{(k)} = \sigma_i^{(k)}e_i^{(k)}, \qquad \sigma_i^{(k)}>0.
\end{eqnarray*}
Taking the union of these bases over all $k$ gives an orthonormal
basis of $\mathcal H$. Relative to this basis every $S_k$ is
diagonal, and consequently $S$ is block diagonal. Furthermore,
$$T(\mathcal H_k)\subseteq\mathcal H_{k+1},$$ implies that, relative
to the decomposition $\mathcal H=\bigoplus_{k=1}^{\infty}\mathcal
H_k$, the operator $T$ has the block representation
\begin{eqnarray*}
T=
\begin{bmatrix} 0&0&0&\cdots\\ T_1&0&0&\cdots\\ 0&T_2&0&\cdots\\
\vdots&\vdots&\vdots&\ddots \end{bmatrix},
\end{eqnarray*}
where, $T_k:=T|_{\mathcal H_k}:\mathcal H_k \rightarrow \mathcal
H_{k+1}$. Thus $T$ is strictly lower triangular. This proves
$(i)$. To prove $(ii)$, define $A:=S-T$. Since
$\|S^{-1}T\|=\rho<1$, the Neumann series theorem implies that
$I-S^{-1}T$ is invertible. Because $A = S(I-S^{-1}T)$, and $S$ is
invertible as the frame operator of an $AR$-frame, it follows that
$A=S-T$ is invertible. The iteration (\ref{eqn-90}) can be
rewritten as
\begin{eqnarray} \label{eqn-1000} \notag
u_{n+1} &= u_n+S^{-1}(x-Au_n)\\ &= S^{-1}x+S^{-1}Tu_n.
\end{eqnarray}

Let $u=A^{-1}x=(S-T)^{-1}x$. Then $Au=Su-Tu=x$, or
\begin{eqnarray} \label{eqn-1001}
u=S^{-1}x+S^{-1}Tu.
\end{eqnarray}
Define the error $e_n=u-u_n$. Then two identities (\ref{eqn-1000})
and (\ref{eqn-1001}) gives
\begin{eqnarray*}
\begin{aligned} e_{n+1} &=
S^{-1}Tu-S^{-1}Tu_n\\ &= S^{-1}T(u-u_n)\\ &= S^{-1}Te_n.
\end{aligned}
\end{eqnarray*}
Hence, $e_n=(S^{-1}T)^ne_0$. Therefore,
\begin{eqnarray*}
\|e_n\| \le \|S^{-1}T\|^n\|e_0\| = \rho^n\|u\|,
\end{eqnarray*}
because $e_0=u-u_0=u$. Since $0<\rho<1$, we have
$\rho^n\longrightarrow0$, and consequently
\begin{eqnarray*}
u_n\longrightarrow u=(S-T)^{-1}x,
\end{eqnarray*}
in norm. This completes the proof.
\end{proof}
\begin{rem}
This proposition uses $T$ in an essential way. It is an
application of the Richardson (or Picard) iteration
\cite{Richard,Yousef} to the block operator equation $(S-T)u=x$.
The hypothesis $\rho=\|S^{-1}T\|<1$ guarantees convergence by the
Neumann series theorem, while the block structure induced by the
$AR$-frame gives the adaptive multilevel interpretation of the
iteration.
\end{rem}

Our second result establishes the connection between adaptive
multiresolution frame decompositions and level-wise
preconditioning, and identifies the optimal block-diagonal scaling
that minimizes the worst-case convergence factor of the iterative
reconstruction.


\begin{prop} \label{Prop.22}
Let $0<A_k I_{\mathcal H_k} \leq S_k \leq B_k I_{\mathcal H_k},
\,\ k\geq 1$, where $0<A_k\leq B_k<\infty $. Also, let
\begin{eqnarray} \label{eqn-91}
P_{\alpha} = \bigoplus_{k=1}^{\infty}\alpha_k S_k^{-1},
\end{eqnarray}
where $\alpha_k>0$ are scaling parameters satisfy
$\sup_{k\ge1}\frac{\alpha_k}{A_k}<\infty$. Then the following
statements hold:
\begin{itemize}
\item[(i)] $P_{\alpha}\in\mathcal B(\mathcal H)$, and $P_\alpha S
= \bigoplus_{k=1}^{\infty}\alpha_k I_{\mathcal H_k}$. If, in
addition $$0<\alpha_- \leq \alpha_k\leq \alpha_+<\infty ,$$ then
$\alpha_- I \leq P_\alpha S \leq \alpha_+ I$. \item[(ii)] Consider
the iterative reconstruction
\begin{eqnarray} \label{eqn-92}
x_{n+1} = x_n+ P_\alpha(y-Sx_n),
\end{eqnarray}
where $y=Sx$. Then the error operator is $E_\alpha = I-P_\alpha
S$, and the error satisfies $e_{n+1} = E_\alpha e_n$, where
$e_n=x-x_n$. Consequently,
$$\|x-x_n\| \leq \|E_\alpha\|^n\|x-x_0\|.$$
Among all block-diagonal preconditioners of the form
(\ref{eqn-91}), the optimal scalar normalization on each level is
\begin{eqnarray*}
\alpha_k^{*} = \frac{2} {\lambda_{\max}(S_k)+\lambda_{\min}(S_k)}=
\frac{2}{A_k+B_k},
\end{eqnarray*}
for which
\begin{eqnarray*}
\|I-\alpha_k^{*}S_k\| = \frac{B_k-A_k}{B_k+A_k},
\end{eqnarray*}
Consequently,
\begin{eqnarray*}
\|E_{\alpha^*}\| = \sup_{k\geq1} \frac{B_k-A_k}{B_k+A_k}.
\end{eqnarray*}
and this is the smallest possible operator norm among all
block-diagonal preconditioners of the above form.
\end{itemize}
\end{prop}

\begin{proof}
Since $\mathcal H=\bigoplus_{k=1}^{\infty}\mathcal H_k$, every
vector $x\in\mathcal H$ admits the unique orthogonal decomposition
$x=\sum_{k=1}^{\infty}x_k, \,\ x_k\in\mathcal H_k$, with $\|x\|^2
= \sum_{k=1}^{\infty}\|x_k\|^2$. Moreover, $$S=
\bigoplus_{k=1}^{\infty}S_k.$$ Since $A_kI_{\mathcal H_k} \le
S_k$, each $S_k$ is positive and invertible on $\mathcal H_k$.
Furthermore, $\|S_k^{-1}\| = \frac1{\lambda_{\min}(S_k)} \le
\frac1{A_k}$. Hence, for every $x=\sum_kx_k, \,\ P_\alpha x =
\sum_{k=1}^{\infty} \alpha_kS_k^{-1}x_k$, and therefore
\begin{eqnarray*}
\|P_\alpha x\|^2 &=& \sum_{k=1}^{\infty} \alpha_k^2
\|S_k^{-1}x_k\|^2 \\ &\le& \sum_{k=1}^{\infty}
\frac{\alpha_k^2}{A_k^2} \|x_k\|^2 \\ &\le& \left( \sup_{k\ge1}
\frac{\alpha_k}{A_k} \right)^2 \sum_{k=1}^{\infty}\|x_k\|^2.
\end{eqnarray*}
Thus $\|P_\alpha x\| \le \left( \sup_{k\ge1} \frac{\alpha_k}{A_k}
\right)\|x\|$, showing that $P_\alpha\in\mathcal B(\mathcal H)$.
Next,
\begin{eqnarray} \label{eqn-93}
P_\alpha S = \left( \bigoplus_{k=1}^{\infty} \alpha_kS_k^{-1}
\right) \left( \bigoplus_{k=1}^{\infty}S_k \right).
\end{eqnarray}
Since products of block-diagonal operators are computed blockwise,
(\ref{eqn-93}) gives
\begin{eqnarray*}
P_\alpha S = \bigoplus_{k=1}^{\infty}
\alpha_kS_k^{-1}S_k=\bigoplus_{k=1}^{\infty} \alpha_kI_{\mathcal
H_k}.
\end{eqnarray*}
If $0<\alpha_-\le\alpha_k\le\alpha_+$, then for every
$x=\sum_kx_k$,

\begin{eqnarray*}
\langle P_\alpha Sx,x\rangle &= \sum_{k=1}^{\infty}
\alpha_k\|x_k\|^2.
\end{eqnarray*}
Therefore, $\alpha_-\|x\|^2 \le \langle P_\alpha Sx,x\rangle \le
\alpha_+\|x\|^2$, which is equivalent to $$\alpha_-I \le P_\alpha
S \le \alpha_+I.$$ This proves $(i)$.\\
For part $(ii)$, let $e_n=x-x_n $. The iteration (\ref{eqn-92})
gives
\begin{eqnarray*}
e_{n+1} &=& x-x_n-P_\alpha(y-Sx_n)\\ &=& x-x_n-P_\alpha(Sx-Sx_n)\\
&=& (I-P_\alpha S)e_n .
\end{eqnarray*}
Hence, $e_n=(I-P_\alpha S)^ne_0$, and therefore, $\|e_n\| \leq
\|I-P_\alpha S\|^n \|e_0\|$. To determine the optimal
normalization, consider a fixed level $k$. Since $S_k$ is
self-adjoint and positive, the spectrum of $S_k$ satisfies
\begin{eqnarray*}
\sigma(S_k) \subset[A_k,B_k].
\end{eqnarray*}
The restriction of the error operator to $\mathcal H_k$ is
$I-\alpha_kS_k$ whose operator norm satisfies
\begin{eqnarray*}
\|I-\alpha_kS_k\| = \max_{\lambda\in[A_k,B_k]}
|1-\alpha_k\lambda|.
\end{eqnarray*}
The optimal value of $\alpha_k$ minimizes this maximum. Indeed,
since the function $$\lambda \longmapsto |1-\alpha_k\lambda|,$$ is
convex on $[A_k,B_k]$, a basic property of convex functions is
that their maximum over a compact interval is attained at an
endpoint, so its maximum over the interval is attained at one of
the endpoints. Hence
\begin{eqnarray*}
\max_{\lambda\in[A_k,B_k]}|1-\alpha_k\lambda| =
\max\{|1-\alpha_kA_k|,\;|1-\alpha_kB_k|\}.
\end{eqnarray*}
The minimax value is achieved when these two endpoint errors are
equal; otherwise, one can perturb $\alpha_k$ to decrease the
larger error without increasing the smaller one beyond it. Thus
$|1-\alpha_kA_k| = |1-\alpha_kB_k|$, which, since
$1/B_k<\alpha_k<1/A_k$, reduces to $$1-\alpha_kA_k =
\alpha_kB_k-1.$$ Indeed, since $A_k,B_k>0$, the optimal $\alpha_k$
lies between $\frac{1}{B_k}$ and $\frac{1}{A_k}$. Indeed, if
$\alpha_k<1/B_k$, then both $1-\alpha_k A_k$ and $1-\alpha_k B_k$
are positive, and increasing $\alpha_k$ decreases both, so this
cannot be optimal; if $\alpha_k>1/A_k$, then both are negative,
and decreasing $\alpha_k$ decreases the maximum. Thus the
minimizer satisfies $\frac1{B_k} < \alpha_k < \frac1{A_k}$, so
that $1-\alpha_k A_k>0, \,\ 1-\alpha_k B_k<0$. Hence

$$|1-\alpha_k A_k| = 1-\alpha_k A_k, \,\ \mbox{and} \,\ |1-\alpha_k B_k| =
\alpha_k B_k-1.$$ Therefore the equality of endpoint errors
becomes $1-\alpha_k A_k = \alpha_k B_k-1$. Solving for $\alpha_k$
gives
\begin{eqnarray*}
\alpha_k^* = \frac{2}{A_k+B_k}.
\end{eqnarray*}
 Substituting this value, gives
\begin{eqnarray*}
\|I-\alpha_k^*S_k\| &=& 1-\alpha_k^*A_k\\ &=&
1-\frac{2A_k}{A_k+B_k}\\ &=& \frac{B_k-A_k}{B_k+A_k}.
\end{eqnarray*}
Finally,
\begin{eqnarray*}
I-P_{\alpha^*}S = \bigoplus_{k=1}^{\infty} (I-\alpha_k^*S_k),
\end{eqnarray*}
Since the norm of a block-diagonal operator acting on an
orthogonal direct sum equals the supremum of the norms of its
diagonal blocks,
 $\|I-P_{\alpha^*}S\| = \sup_k \|I-\alpha_k^*S_k\|$.
Hence,
\begin{eqnarray*}
\|I-P_{\alpha^*}S\| = \sup_k \|I-\alpha_k^*S_k\| = \sup_k
\frac{B_k-A_k}{B_k+A_k}.
\end{eqnarray*}
Because the minimization problem is independent on each orthogonal
level, no other choice of positive scalars $\{\alpha_k\}$ can
produce a smaller block norm. Hence $P_{\alpha^*}$ is the optimal
block-diagonal preconditioner within the class
\begin{eqnarray*}
\left\{ \bigoplus_{k=1}^{\infty}\alpha_kS_k^{-1}: \alpha_k>0,\,
\sup_k\frac{\alpha_k}{A_k}<\infty \right\},
\end{eqnarray*}
and the proof is complete.

\end{proof}


\begin{thebibliography}{99}
\bibitem[Richard]{Richard} Varga RS. Matrix Iterative Analysis, 2nd ed., Springer, 2000
\bibitem[Yousef]{Yousef} Saad Y. Iterative Methods for Sparse Linear Systems, 2nd ed., SIAM, 2003.
\end{thebibliography}
\end{document}
